%\section{Arquitetura}
%
%\subsection{\textit{Timer}}
%
%O componente \textit{Timer} é responsável pelo gerenciamento temporal do sistema, fornecendo mecanismos de controle e sincronização para eventos periódicos.
%Sua principal função consiste em atualizar e notificar os módulos dependentes a cada interrupção temporal (\textit{tick}).
%Além disso, esse componente disponibiliza interfaces para consulta do tempo corrente, permitindo que outros módulos realizem operações dependentes de temporização de forma padronizada.
%
%\vspace{0.5cm}
%
%\begin{tikzpicture}
%\umlclass[
%fill=white,
%draw=black,
%text=black
%]{Timer}{
%- tickers : \texttt{Tuple<Tickers...>}
%}{
%+ init() : \texttt{void} \\
%+ onTick() : \texttt{void} \\
%+ now() : \texttt{Microsecond} \\
%+ us() : \texttt{Microsecond} \\
%}
%\end{tikzpicture}
%
%\subsection{\textit{VirtualCPU}}
%
%%O componente \textit{Timer} é responsável pelo gerenciamento temporal do sistema, fornecendo mecanismos de controle e sincronização para eventos periódicos.
%
%\vspace{0.5cm}
%
%\begin{tikzpicture}
%\umlclass[
%fill=white,
%draw=black,
%text=black
%]{VirtualCPU}{
%- vm : \texttt{VirtualMachine*}
%}{
%+ activate() : \texttt{void} \\
%+ interrupt(id : \texttt{unsigned int}) : \texttt{void} \\
%+ onTick() : \texttt{void} \\
%+ read() : \texttt{bool} \\
%+ write() : \texttt{bool} \\
%}
%\end{tikzpicture}
